% answers.tex —— 1958 年·湖南卷 · 参考答案与详细解析片段
% 由 tools/generate_projects.py 依来源总账自 1958/1958_湖南_数学真题_试题.md + 1958/1958_湖南_数学真题_解析正文.md + 1958/1958_湖南_数学真题_答案要点.md 生成，勿手改（改源 MD 后重跑）


\begin{answer}
\gksolhead{1}{}

\textbf{题目}：

不用除法证明 $(2+x)^{2n-1}+x^{2n-1}+x+1$ 能被 $x+1$ 整除（$n$ 为正整数）。

\Needspace{4\baselineskip}
\textbf{解析}

\Needspace{5\baselineskip}
设 $f(x)=(2+x)^{2n-1}+x^{2n-1}+x+1$。多项式能被 $x+1$ 整除的充要条件是其在 $x=-1$ 处取值为零（余数定理），于是不必做任何除法，只需计算 $f(-1)$。由 $n$ 为正整数知 $2n-1$ 为奇数，故 $(-1)^{2n-1}=-1$，且 $2+x$ 在 $x=-1$ 时为 $1$：

\nopagebreak[4]
\[
f(-1)=(2-1)^{2n-1}+(-1)^{2n-1}+(-1)+1=1-1-1+1=0
\]

由余数定理，$x+1$ 整除 $f(x)$，即 $(2+x)^{2n-1}+x^{2n-1}+x+1$ 能被 $x+1$ 整除。

\gkanswer{答案：见解析版}

\end{answer}
\begin{answer}
\gksolhead{2}{}

\textbf{题目}：

不解方程 $6x^2-x-2=0$，求作一个二次方程，使它的根是原方程的根的倒数。

\Needspace{4\baselineskip}
\textbf{解析}

\Needspace{5\baselineskip}
设原方程 $6x^2-x-2=0$ 的两根为 $x_1,x_2$，由韦达定理

\nopagebreak[4]
\[
x_1+x_2=\frac{1}{6},\qquad x_1x_2=-\frac{2}{6}=-\frac{1}{3}
\]

\Needspace{5\baselineskip}
所求方程的根为 $\dfrac{1}{x_1},\dfrac{1}{x_2}$，其和与积为

\nopagebreak[4]
\[
\frac{1}{x_1}+\frac{1}{x_2}=\frac{x_1+x_2}{x_1x_2}=\frac{1/6}{-1/3}=-\frac{1}{2},\qquad
\frac{1}{x_1}\cdot\frac{1}{x_2}=\frac{1}{x_1x_2}=-3
\]

\Needspace{5\baselineskip}
故所求二次方程为

\nopagebreak[4]
\[
y^2-\left(-\frac12\right)y+(-3)=0,\qquad\text{即}\ y^2+\frac12y-3=0
\]

两边乘 $2$ 得 $2y^2+y-6=0$（未知数记号可任取，即 $2x^2+x-6=0$）。

\gkanswer{答案：$2x^2+x-6=0$}

\end{answer}
\begin{answer}
\gksolhead{3}{}

\textbf{题目}：

求证 $\sin 75^\circ=\dfrac14\left(\sqrt6+\sqrt2\right)$。

\Needspace{4\baselineskip}
\textbf{解析}

\Needspace{5\baselineskip}
把 $75^\circ$ 拆成两个特殊角之和，用两角和的正弦公式：

\nopagebreak[4]
\[
\sin 75^\circ=\sin(45^\circ+30^\circ)=\sin45^\circ\cos30^\circ+\cos45^\circ\sin30^\circ
\]

\Needspace{5\baselineskip}
代入 $\sin45^\circ=\cos45^\circ=\dfrac{\sqrt2}{2}$，$\sin30^\circ=\dfrac12$，$\cos30^\circ=\dfrac{\sqrt3}{2}$：

\nopagebreak[4]
\[
\sin 75^\circ=\frac{\sqrt2}{2}\cdot\frac{\sqrt3}{2}+\frac{\sqrt2}{2}\cdot\frac12=\frac{\sqrt6}{4}+\frac{\sqrt2}{4}=\frac{\sqrt6+\sqrt2}{4}
\]

\gkanswer{答案：见解析版}

\end{answer}
\begin{answer}
\gksolhead{4}{}

\textbf{题目}：

证明：设三角形的三个角成等差数列，最大角与最小角相差 $30^\circ$，则其三边之比顺次等于 $2\sqrt2:2\sqrt3:(\sqrt6+\sqrt2)$ 之比。

\Needspace{4\baselineskip}
\textbf{解析}

\Needspace{5\baselineskip}
设三个内角为 $\beta-d,\ \beta,\ \beta+d$（等差数列，公差 $d$）。由内角和定理

\nopagebreak[4]
\[
(\beta-d)+\beta+(\beta+d)=3\beta=180^\circ,\qquad \beta=60^\circ
\]

\Needspace{5\baselineskip}
最大角与最小角相差 $(\beta+d)-(\beta-d)=2d=30^\circ$，得 $d=15^\circ$。故三个角为

\nopagebreak[4]
\[
45^\circ,\quad 60^\circ,\quad 75^\circ
\]

\Needspace{5\baselineskip}
由正弦定理，三边之比等于其对角正弦之比（设 $a,b,c$ 分别为 $45^\circ,60^\circ,75^\circ$ 的对边）：

\nopagebreak[4]
\[
a:b:c=\sin45^\circ:\sin60^\circ:\sin75^\circ
\]

\Needspace{5\baselineskip}
其中 $\sin45^\circ=\dfrac{\sqrt2}{2}$，$\sin60^\circ=\dfrac{\sqrt3}{2}$，$\sin75^\circ=\dfrac{\sqrt6+\sqrt2}{4}$（见第 3 题）。三数统一乘 $4$：

\nopagebreak[4]
\[
a:b:c=\frac{\sqrt2}{2}:\frac{\sqrt3}{2}:\frac{\sqrt6+\sqrt2}{4}=2\sqrt2:2\sqrt3:(\sqrt6+\sqrt2)
\]

即三边之比顺次（从小角的对边到大角的对边）等于 $2\sqrt2:2\sqrt3:(\sqrt6+\sqrt2)$。

\gkanswer{答案：见解析版}

\end{answer}
\begin{answer}
\gksolhead{5}{}

\textbf{题目}：

一矩形内接于圆，它的面积是 $48\,\text{cm}^2$，圆半径是 $5\,\text{cm}$，求矩形的两邻边。

\Needspace{4\baselineskip}
\textbf{解析}

\Needspace{5\baselineskip}
矩形内接于圆，其对角线恰为圆的直径，故对角线长

\nopagebreak[4]
\[
2\times5=10
\]

\Needspace{5\baselineskip}
设两邻边为 $a,b$（cm）。由矩形面积与勾股定理：

\nopagebreak[4]
\[
ab=48,\qquad a^2+b^2=10^2=100
\]

\Needspace{5\baselineskip}
于是

\nopagebreak[4]
\[
(a+b)^2=a^2+b^2+2ab=100+96=196,\qquad a+b=14
\]

\Needspace{5\baselineskip}
$a,b$ 是二次方程 $t^2-14t+48=0$ 的两根：

\nopagebreak[4]
\[
t=\frac{14\pm\sqrt{14^2-4\cdot48}}{2}=\frac{14\pm2}{2}
\]

得 $t=6$ 或 $t=8$。故矩形的两邻边为 $6\,\text{cm}$ 与 $8\,\text{cm}$。

\gkanswer{答案：$6\,\text{cm}$ 与 $8\,\text{cm}$}

\end{answer}
\begin{answer}
\gksolhead{6}{}

\textbf{题目}：

设 $AD,BE$ 是由圆的直径 $AB$ 两端点所引的切线，$DE$ 是过圆上任意一点 $F$ 的切线，此切线与 $AD,BE$ 相交于 $D$ 及 $E$，求证：$AD\cdot BE=OF^2$。

\Needspace{4\baselineskip}
\textbf{解析}

（记号说明：$O$ 为圆心，原卷按通行记号未再显式定义。）

由切线长定理：自圆外一点所引的两条切线段相等，且圆心与该点的连线平分两切线的夹角。于是

\begin{itemize}
\item $D$ 处：$DA=DF$，且 $OD$ 平分 $\angle ADE$；
\item $E$ 处：$EB=EF$，且 $OE$ 平分 $\angle DEB$。
\end{itemize}

\Needspace{5\baselineskip}
又切线垂直于过切点的半径，故 $AD\perp AB$，$BE\perp AB$，从而 $AD\parallel BE$，同旁内角互补：

\nopagebreak[4]
\[
\angle ADE+\angle DEB=180^\circ
\]

\Needspace{5\baselineskip}
两角各被平分一半：

\nopagebreak[4]
\[
\angle DOE=180^\circ-(\angle ODE+\angle OED)=180^\circ-\frac{180^\circ}{2}=90^\circ
\]

\Needspace{5\baselineskip}
即 $\triangle DOE$ 是以 $DE$ 为斜边的直角三角形。而 $OF\perp DE$（半径垂直于过切点 $F$ 的切线），即 $OF$ 是直角三角形斜边上的高。由射影定理（斜边上的高是两段分线的比例中项）：

\nopagebreak[4]
\[
OF^2=DF\cdot FE
\]

\Needspace{5\baselineskip}
代入切线长 $DF=DA$、$FE=EB$，得

\nopagebreak[4]
\[
OF^2=AD\cdot BE
\]

\gkanswer{答案：见解析版}

\end{answer}
\begin{answer}
\gksolhead{7}{}

\textbf{题目}：

正六棱锥底面的边长是 $4\,\text{cm}$，侧面与底面所成二面角等于 $45^\circ$，求这棱锥体的体积。

\Needspace{4\baselineskip}
\textbf{解析}

\Needspace{5\baselineskip}
底面是边长 $4\,\text{cm}$ 的正六边形，其面积为

\nopagebreak[4]
\[
S_{\text{底}}=6\cdot\frac{\sqrt3}{4}\cdot4^2=24\sqrt3\ (\text{cm}^2)
\]

\Needspace{5\baselineskip}
底面的边心距（中心到边的距离）为

\nopagebreak[4]
\[
r=4\cdot\frac{\sqrt3}{2}=2\sqrt3\ (\text{cm})
\]

\Needspace{5\baselineskip}
正棱锥的高、底面边心距与侧面斜高共处于一个直角三角形中，侧面与底面所成二面角的平面角就是斜高与边心距的夹角，即 $45^\circ$。于是高

\nopagebreak[4]
\[
h=r\tan45^\circ=2\sqrt3\times1=2\sqrt3\ (\text{cm})
\]

\Needspace{5\baselineskip}
体积为

\nopagebreak[4]
\[
V=\frac13\,S_{\text{底}}\cdot h=\frac13\cdot24\sqrt3\cdot2\sqrt3=\frac13\cdot144=48\ (\text{cm}^3)
\]

\gkanswer{答案：$48\,\text{cm}^3$}

\end{answer}
\begin{answer}
\gksolhead{8}{}

\textbf{题目}：

证明二面角的平分面内任意一点到这二面角的两个面的距离相等。

\Needspace{4\baselineskip}
\textbf{解析}

设二面角的两个面为 $\alpha,\beta$，棱为 $l$，平分面为 $\gamma$（$\gamma$ 过棱 $l$），$P$ 为 $\gamma$ 内任意一点。

过 $P$ 作平面 $\pi\perp l$，设 $\pi$ 交棱 $l$ 于 $Q$，交面 $\alpha$ 于射线 $QA$，交面 $\beta$ 于射线 $QB$，交平分面 $\gamma$ 于射线 $QP$。

因 $l\perp\pi$ 且 $l\subset\alpha,\ l\subset\beta,\ l\subset\gamma$，故 $\alpha,\beta,\gamma$ 都垂直于 $\pi$，而二面角 $\alpha\text{-}l\text{-}\beta$ 的平面角就是 $\pi$ 内的 $\angle AQB$；由 $\gamma$ 平分二面角，射线 $QP$ 平分 $\angle AQB$。

在平面 $\pi$ 内，过 $P$ 作 $PA'\perp QA$（垂足 $A'$），$PB'\perp QB$（垂足 $B'$）。由 $\alpha\perp\pi$，$\pi$ 内垂直于交线 $QA$ 的直线必垂直于面 $\alpha$，故 $PA'\perp\alpha$，$PA'$ 的长就是 $P$ 到面 $\alpha$ 的距离；同理 $PB'\perp\beta$，$PB'$ 的长就是 $P$ 到面 $\beta$ 的距离。

\Needspace{5\baselineskip}
在 $\triangle PA'Q$ 与 $\triangle PB'Q$ 中：$\angle PA'Q=\angle PB'Q=90^\circ$，$PQ$ 为公共边，$\angle PQA'=\angle PQB'$（$QP$ 平分 $\angle AQB$），故两三角形全等（角角边），得

\nopagebreak[4]
\[
PA'=PB'
\]

即 $P$ 到两个面的距离相等。由 $P$ 在 $\gamma$ 内的任意性，二面角平分面内任意一点到这二面角的两个面的距离相等。

\gkanswer{答案：见解析版}

\end{answer}